\section{The System Through the Multi-Agent Paradigm}
\label{sec:mas}

The preceding sections described \emph{what} was built and \emph{how} it is
trained. This section changes register and asks what kind of system it is.
Multi-agent systems are not merely a technology but an engineering perspective:
a way of decomposing a problem into autonomous, interacting, situated components
and of locating where a system's intelligence actually resides
\citep{jennings2001agent,wooldridge2009introduction}. Read through that lens the
drone team instantiates a remarkably wide slice of the multi-agent conceptual
apparatus, from agenthood and the environment-as-abstraction through
coordination, artefacts, self-organization, resilience and graduated autonomy to
multi-agent learning. Throughout, every conceptual claim is anchored to a
concrete mechanism from Sections~\ref{sec:problem}--\ref{sec:training}, and the
section closes by asking, honestly, what a richer multi-agent design would add
and at what cost.

\subsection{Why a multi-agent formulation}
\label{subsec:mas-why}

It is worth beginning with the question the rest of the section presupposes: why
is this a multi-agent system at all, rather than one program steering four
vehicles? The answer is that the defining constraints of the task, not
implementation taste, rule the centralized alternative out along the three
dimensions that classically motivate the agent paradigm: distribution, autonomy
and interaction \citep{jennings2001agent,zambonelli2003signs}.

\emph{Distribution of sensing and knowledge.} No component of the system ever
holds the global state. Each drone perceives a Chebyshev window of radius
$\rvis$ and knows, beyond that, only what it has accumulated or been told. A
central controller would need the union of all observations at every turn, but
the communication model is range-limited: drones farther apart than $\rcomm$
cannot exchange anything. Centralized control over a communication-limited team
is not merely inefficient; during the frequent, deliberate intervals in which the
team is partitioned into disconnected communication components, it is
\emph{physically unimplementable}. The problem itself is distributed, so the
control must be.

\emph{Autonomy as a fault-tolerance requirement.} The fault model sharpens the
argument to its most concrete form. A centralized controller is a single point of
failure, and this project's central premise is that components \emph{do} fail, at
random, mid-mission. If the deciding entity resided in one drone and that drone
broke, the team would be decapitated. In the actual design, every drone carries
the full decision-making apparatus, its own copy of the shared policy, and the
death of any subset of drones leaves every survivor as capable as before.
Autonomy here is not an aesthetic commitment: it is what makes graceful
degradation possible at all.

\emph{Interaction as the source of system-level competence.} Finally, the task
is one no single drone solves well: a lone searcher covers area linearly in time,
while a well-dispersed team covers it in parallel. The system-level competence,
sweeping the map quickly, avoiding duplicated effort and re-covering the sectors
of fallen teammates, is located in no individual policy evaluation; it arises
\emph{between} the drones, from their interactions through fusion, mutual
observation and shared spatial traces. That is precisely what distinguishes a
multi-agent system from a parallel implementation of a single-agent one
\citep{wooldridge2009introduction}.

One clarification prevents a recurring misunderstanding: the drones share network
\emph{weights}, and weight sharing is a training-time device, not a runtime
coupling. At execution each drone evaluates its own copy of the function on its
own inputs; no activation, gradient or decision crosses drone boundaries. Four
copies of the same brain, fed four different experiences, are four agents, just
as identical twins are two people.

\subsection{Drones as agents}
\label{subsec:mas-agents}

``Agent'' is a contested term, but a broadly shared core defines it as a
computational system situated in an environment and capable of autonomous action
in pursuit of its design objectives, with the \emph{weak notion of agency}
requiring autonomy, reactivity, proactiveness and social ability
\citep{wooldridge1995intelligent,franklin1996agent}. This subsection audits the
drone against each requirement, and is equally explicit about which stronger
notions the drone does \emph{not} satisfy.

\paragraph{Autonomy: encapsulated control.}
The most fundamental computational reading of autonomy is encapsulation of
control: an agent is a locus of decision that cannot be \emph{invoked}, other
components being able to influence it only through data, never by calling its
behaviour as one calls a procedure \citep{wooldridge1995intelligent}. The drones
satisfy this literally. Nothing in the system commands a drone: teammates cannot
move it, interrupt it or query it, and the only thing that ever crosses the
boundary between two drones is world knowledge, through fusion. Each turn
consists of reading \emph{its own} observation and context, evaluating \emph{its
own} policy copy, and acting. Even the round-robin scheduling does not breach
encapsulation: it serializes \emph{when} drones act, a property of the
environment's clock, not \emph{what} they decide.

\paragraph{Situatedness.}
Agents act \emph{in} and \emph{through} an environment, and their action cannot
be divorced from the context in which it occurs \citep{suchman1987plans}. The
drone's entire interface to the world is environmental: it perceives by revealing
cells around its position, acts by moving to adjacent cells, and even its
communication is situated, fusion being triggered by spatial proximity rather
than by addressing a recipient. There is no action-at-a-distance anywhere in the
system, and the policy's inputs are spatial through and through: position, fault
knowledge and history all reach the network as \emph{maps}, an architectural echo
of the fact that for this agent context is location.

\paragraph{Reactivity, with state.}
The classical formalization of a non-trivially reactive agent equips it with
internal state: a perception function from environment states to percepts, a
state-update function folding each percept into an internal state, and an action
function from internal states to actions \citep{wooldridge2009introduction}. The
drone realizes the triple exactly: perception is the vision update (cells within
$V_i^t$, plus beacon detection); the state update is the monotone accumulation
into $K_i^t$, merged with teammates' contributions by fusion; the action function
is the Q-network's masked greedy selection over the rendering of $K_i^t$. That
internal state is what lifts the drone above a thermostat-style pure reflex: two
drones on the same cell with the same vision act differently if their
\emph{histories} differ, because their maps do.

\paragraph{Proactiveness: distilled, not represented.}
Agents are expected to exhibit goal-directed behaviour, taking initiative rather
than merely responding to stimuli \citep{wooldridge1995intelligent}. The drone is
unambiguously goal-directed in behaviour, since it systematically expands
coverage, avoids re-searching and homes in on the target once its location enters
the map, yet nowhere in the drone does the goal \emph{exist as an object}. There
is no symbol denoting ``find the target'', no explicit desire, no plan. The
goal-directedness was distilled at training time from the reward function of
Table~\ref{tab:reward} into the value estimates that now drive action selection.
In the classical dichotomy, the drone is \emph{goal-oriented} rather than
\emph{goal-governed}: its objective is implicit in the structure of its behaviour
rather than explicitly represented and reasoned over. This is a genuine
trade-off, not a defect, and Section~\ref{subsec:mas-beyond} discusses what
explicit goals would buy, and cost.

\paragraph{Social ability.}
The drone's social repertoire is narrow but real: it donates its accumulated
knowledge to whoever comes within $\rcomm$, incorporates teammates' knowledge
symmetrically, perceives nearby live teammates through the Other\_Position
channel, and both emits, passively as a wreck, and consumes distress information.
The next subsection analyses why this thin, language-free interaction layer is
nonetheless a coordination mechanism of some sophistication.

\paragraph{What the drone is not.}
Against the \emph{strong} notion of agency, with mental states, beliefs and
intentions in the technical sense of belief--desire--intention architectures
\citep{rao1995bdi,bratman1987intention}, the drone plainly falls short, and it is
worth being precise about where. The per-drone knowledge state is
belief-\emph{like}: private, partial, possibly stale, updated by observation and
testimony. But it is a fixed-schema spatial data structure, not a revisable base
of propositions; the drone cannot represent ``teammate 2 probably went north'',
let alone reason about what teammate 2 believes. Deliberation, in the sense of
weighing candidate intentions, is likewise absent, replaced by a single amortized
function evaluation. The system thus occupies a definite point in the agent
design space: fully agentive in the weak sense, deliberately sub-cognitive in the
strong one.

\subsection{The environment as a first-class abstraction}
\label{subsec:mas-environment}

A recurring lesson of multi-agent engineering is that the environment is not
scaffolding around the agents but a design dimension in its own right: an active
entity with its own state, dynamics and rules, through which agents perceive, act
and interact \citep{weyns2007environment}. This project embodies that stance
structurally: \texttt{DroneSearchEnv} is not a thin arena but the locus of most of
the system's semantics. The agent package is larger in absolute size, but almost
all of it is network definition and learning machinery; the world model,
including every rule of interaction, lives entirely on the environment side.
Four environmental
responsibilities deserve explicit notice. \emph{Topology and locality}: the grid,
its Chebyshev vision neighbourhoods and its Manhattan communication distances
define what ``near'' means, and thereby delimit perception, interaction and
fusion; locality is not a property agents happen to have but a rule the
environment enforces uniformly on everyone. \emph{Dynamics}: the fault process
lives in the environment, not in the agents, since drones do not decide to fail,
the world fails them, and the same holds for the episode clock, including the
rule that wrecks' turns still consume budget. \emph{Mediation}: every interaction
between drones, whether fusion, mutual sighting or beacon detection, is computed
by the environment as a side effect of spatial configuration, so drones have no
channel to each other that bypasses the world, which is exactly the
situated-interaction discipline \citet{weyns2007environment} argue for. And
\emph{governance}: the action mask constrains the space of legal behaviour at its
source, while the reward terms of Table~\ref{tab:reward} are environmental
legislation, the incentive structure under which autonomous behaviour is shaped
without ever being commanded.

The methodological payoff appeared already in Section~\ref{sec:architecture}:
because the environment owns the rules, the agents could be swapped (random,
keyboard-driven, BFS-guided, learned) with zero change to the world's semantics,
which is what lets Section~\ref{sec:experiments} run the learned policy and the
BFS auto-navigation override over identical seeded instances.

\subsection{Interaction and coordination}
\label{subsec:mas-coordination}

Coordination, in the sense of managing the interdependencies between concurrent
activities so that a collection of agents behaves as a coherent ensemble, is the
central problem any multi-agent system must solve. This system solves it with a
combination that is instructive precisely because of what it lacks: there are no
messages, no protocols, no negotiation, no roles assigned by any authority.
Coordination arises from three mutually reinforcing mechanisms.

\paragraph{(a) Knowledge-level communication without a language.}
When two drones fuse maps, what travels is world knowledge (``these cells are
free, these are obstacles, the target is here, teammate 3 crashed there'') and
nothing else: no requests, commands, commitments, or any of the performatives
that agent communication languages provide \citep{finin1994kqml}. The omission is
principled rather than lazy. Speech acts exist to coordinate \emph{deliberating}
agents: a request makes sense when the receiver can decide whether to honour it.
Between sub-cognitive policies there is no deliberation to address, and the task
is identical-interest, so there is nothing to negotiate over. What remains
genuinely useful to exchange is exactly what \emph{is} exchanged: factual
knowledge, whose fusion is monotone (element-wise maximum) and therefore
order-independent, idempotent and conflict-free. The design question ``what
should agents say to each other?'' is answered here with ``nothing; they should
compare notes''.

\paragraph{(b) The fused map as a shared dataspace.}
The resulting pattern has the character of a \emph{blackboard} or generative
dataspace rather than of message passing
\citep{nii1986blackboard,gelernter1985generative}: the three properties of
generative communication, namely that emitted information outlives its producer,
is accessible by content rather than by address, and decouples interacting
parties in time and identity, all hold of the fused knowledge. A cell discovered
by drone 1 at round 10 shapes drone 4's decision at round 150, long after the
encounter that transmitted it, possibly via a relay chain (1 fused with 2, later
2 with 4) such that 1 and 4 were \emph{never} within range of each other; and
because fusion is an anonymous maximum, knowledge carries no provenance, the
receiver acting on ``the target is at $(r,c)$'' identically whether it saw the
target itself or inherited the fact third-hand. The team's effective medium of
coordination is thus a distributed, eventually-consistent shared map, differing
from a classical blackboard in having no central board: every drone carries its
own replica, and proximity events are the synchronization protocol.

\paragraph{(c) Stigmergic traces, virtualized.}
The third mechanism is coordination through the \emph{traces} of activity, the
pattern Grass\'e named stigmergy: work performed modifies the environment, and
the modified environment directs subsequent work
\citep{grasse1959reconstruction,parunak1997go}. The Visited and Trajectory
channels are precisely such traces. A drone's movement writes a growing mark into
the shared knowledge (visited cells, own visit counts); those marks then
\emph{repel} future search, economically via the revisit penalty and the
nullified exploration bonus, behaviourally via whatever avoidance the trained
policy has internalized, steering drones away from consumed regions exactly as an
evaporating pheromone steers ants away from depleted paths, with the sign
inverted: repulsion, not attraction. One twist distinguishes this from textbook
stigmergy. The traces are not written into the physical grid, the world's cells
having no memory of visits, but into the drones' \emph{shared beliefs},
propagating through fusion rather than sitting at a location. It is stigmergy
over a distributed cognitive medium rather than over matter; the coordination
logic (work marks the field; marks direct work) is identical.

\paragraph{Learned conventions instead of designed protocols.}
What binds these mechanisms into actual coherence, namely \emph{who} goes where
and \emph{how} the team splits, is not specified anywhere. It is a convention
learned by co-training. During training, all drones' experience flows into one
network; behaviours that reduce duplicated coverage were systematically
reinforced by the team's shared exposure to the collision penalty, the
team-novelty exploration bonus and the shared terminal reward. The identity input
of Section~\ref{subsec:arch-context} gives the convention something to condition
on: from the very first round, when all drones stand on the same cell with
identical maps, the policy can map \emph{different identities to different
headings}, realizing an implicit division of the map without any sector ever
being represented, assigned or agreed upon. This is coordination knowledge stored
in weights rather than in protocol state, with consequences, adaptivity and
opacity in equal measure, that Sections~\ref{subsec:mas-beyond}
and~\ref{subsec:eth-responsibility} examine.

\subsection{An artefact perspective}
\label{subsec:mas-artefacts}

The agents-and-artefacts view of multi-agent systems holds that a MAS is
populated not only by agents but by \emph{artefacts}: non-autonomous,
function-oriented computational entities that agents use, which mediate their
activities, and in which designers can embed coordination machinery
\citep{omicini2008artifacts}. Artefacts do not pursue goals; they provide
operations, behave predictably, and serve. The distinction is illuminating here
because it cleanly separates the system's two kinds of components, and locates
several load-bearing pieces of the design on the artefact side of the line.

The clearest instance is the \emph{accumulated map}: a cognitive artefact with a
definite function (aggregate spatial knowledge), a usage interface (written by
moving and sensing, read at every decision as the network's input) and entirely
reactive, predictable dynamics (monotone accumulation and max-merge; it never
acts on its own). It mediates agent--environment interaction, since the drone
acts on the world \emph{as represented by its map}, and through fusion
agent--agent interaction as well: the fused map is a coordination artefact in
exactly the sense in which tuple spaces and blackboards are
\citep{gelernter1985generative,omicini2008artifacts}. The same reading covers the
\emph{distress beacon}, since a wreck is, in death, reduced to an artefact that
manifests one fact to whoever passes, and the \emph{action mask}, a constraining
artefact that does not decide for the agent but delimits its option space,
artefacts enabling and constraining in the same gesture.

Two observations sharpen the picture. Artefacts are where the \emph{designer's}
intelligence lives at runtime: the policy is learned and opaque, while the map
semantics, the fusion rule, the mask and the reward machinery are designed and
transparent; and the system's verifiable properties (fusion never loses
knowledge, masks never permit illegal moves, wrecks never block corridors) are
all artefact properties, inspectable and predictable by construction
\citep{omicini2008artifacts}. And the tooling of
Section~\ref{subsec:arch-tooling} extends the taxonomy to the humans in the loop:
the GUI, the TensorBoard scalars and the seeded evaluation harness are
observation artefacts through which the \emph{engineer} perceives the MAS, the
instrumentation layer that makes an otherwise opaque adaptive system engineerable
at all.

\subsection{Self-organization and emergence}
\label{subsec:mas-emergence}

Self-organization denotes the acquisition of structure by a system through
internal, distributed processes, without external control
\citep{ashby1962principles,camazine2001self}; emergence denotes the appearance of
coherent macro-level patterns not reducible to, nor explicitly encoded in, the
components' individual specifications \citep{goldstein1999emergence}. Both terms
apply to this system, but sloppily applied they explain nothing; the value lies
in stating precisely \emph{what} is designed and \emph{what} emerges.

\emph{Designed} are the local rules: the reward terms a drone experiences (all
computable from its own transition), the fusion rule, the fault process and the
training procedure itself. \emph{Not designed}, that is nowhere represented in
code or in weights as an explicit specification, are the team-level regularities
the trained system is expected to exhibit: the early-episode radial dispersion
from the shared spawn corner; the partition of the map into implicit,
identity-correlated search sectors; the complementarity of coverage, meaning low
overlap between drones' swept regions; and the re-absorption of a dead drone's
sector by its survivors. Each of these is a macro-pattern over the joint
trajectory of the team, produced by four independent evaluations of one function
that has never been trained on, or even given access to, any team-level
objective, each drone's inputs remaining strictly local throughout.

The classical checklist for self-organized coordination
\citep{camazine2001self}, namely multiple interacting units, local information
only, no external director, positive and negative feedback, is met, with the
mechanisms identifiable one by one: negative feedback is the repulsion from
consumed regions (revisit penalty, spent exploration bonus, collision penalty)
that spreads the team out; positive feedback is the attraction created when the
target's location enters the shared map and propagates through fusion, collapsing
the team's search into a pursuit. The balance of the two is exactly what the
policy has had to learn.

Two honest qualifications distinguish this from natural self-organization. First,
the organization is \emph{trained into} the system: the attractors of the
collective dynamics were shaped by fifty thousand episodes of gradient descent
against a designer-chosen reward, which is self-organization at execution time
and teleology at training time. Second, these team-level regularities have the
status of \emph{design intent with partial quantitative support}: of the three
signatures defined as their empirical test,
Section~\ref{subsec:exp-qualitative} delivers coverage complementarity and leaves
the two trajectory-level ones (dispersion over time, post-fault re-covering
latency) to future instrumentation.

\subsection{Fault tolerance, openness, and resilience}
\label{subsec:mas-resilience}

Multi-agent systems are routinely credited with robustness through redundancy and
decentralization; this project turns that slogan into a concrete, testable
mechanism. Three layers of the design implement it.

\paragraph{Structural redundancy.}
Every drone is a complete, interchangeable searcher: same policy, same sensing
model, same authority (none over the others). The task requires \emph{some} drone
to reach the target, not any particular one, so no individual's loss is
qualitatively worse than another's; combined with the environment-level
guarantees that wrecks neither block corridors nor cost reachability, any single
survivor retains in principle the ability to complete the mission alone, degraded
in time but not in kind. This is the redundancy-based robustness familiar from
biological collectives \citep{camazine2001self}, engineered deliberately.

\paragraph{A weak form of openness.}
The team's composition changes at runtime: an episode that begins with four
agents may end with one. In multi-agent terms the system is \emph{semi-open}:
departures occur unannounced, but no agent ever joins. Even this weak openness
has real consequences. Nothing in any drone's machinery assumes a fixed team size
(the policy is conditioned on team size and believed alive count, and per-agent
structures are allocated per episode), and no drone's correct functioning depends
on any other drone existing. The design would extend to arrivals, the context and
fusion machinery being indifferent to \emph{which} drones exist, though the
trained policy has never experienced them: a gap between the architecture's
openness and the policy's that is easily conflated.

\paragraph{Epistemic decentralization: acting on belief, not knowledge.}
The most distinctive layer is informational. No drone is told the team's true
state: each maintains $\mathcal{C}_i^t$, the crashed teammates it knows of, and
acts on the belief $\hat n_i^t = \nagents - |\mathcal{C}_i^t|$. This is a
genuine, if minimal, epistemic structure: a private, first-order factual belief
state, updated by two canonically distinct sources, direct observation (entering
beacon range or sight of a crash) and testimony (fusion with a better-informed
teammate), with the propagation dynamics of a gossip process, fault knowledge
spreading through the chance encounters of the communication graph at a rate
governed by $\rcomm$ and by how the policy moves the drones. The soundness
property established earlier, that $\hat n_i$ never undercounts the living, has a
behavioural face: a drone's error is always in the direction of
\emph{overestimating} its team, continuing to rely, for a while, on a searcher
that no longer exists. The interval between a fault and a given survivor's
discovery of it is thus the system's window of miscoordination, and its cost is
an empirical question that Section~\ref{subsec:exp-faults} takes up. The design
decision worth defending here is the refusal of the tempting alternative, a
global fault oracle broadcasting team state, which would be simpler and would
quietly reintroduce exactly the centralized, distance-insensitive information
channel whose absence defines the problem.

Note finally how the three layers compose: redundancy makes survivors
\emph{sufficient}, epistemic decentralization makes them eventually \emph{aware},
and context conditioning makes them \emph{adaptive}, since $\hat n_i$ feeds the
policy, so discovering a death is not merely bookkeeping but an input that
licenses recalibrating the division of labour. Fault tolerance is not a module in
this system; it is a property of the composition.

\subsection{The autonomy spectrum}
\label{subsec:mas-autonomy}

Autonomy is graduated, not binary, a lesson institutionalized in the level
taxonomies of applied domains, from driving automation's six levels
\citep{sae2021j3016} to the human-delegation hierarchies of military doctrine
\citep{dod2012autonomy}. Placing this system honestly on that spectrum matters
twice: once for conceptual hygiene, and once because the placement determines
which ethical questions Section~\ref{sec:ethics} must answer.

The distinction that organizes the spectrum is between \emph{automatic} systems,
executing pre-specified stimulus-response rules however elaborate, and
\emph{autonomous} ones, which select and adapt their own courses of action toward
a goal under conditions their designers did not enumerate. The drones sit
unambiguously on the autonomous side: no trajectory, sweep pattern or contingency
table exists anywhere in the system, the mapping from situations to actions was
learned, and it adapts online to circumstances (team size, sensing regime,
obstacle layout, teammate deaths) that vary per episode and per turn. The BFS
auto-navigation override is the clarifying foil: \emph{that} is an automatic
component, a fixed and verifiable procedure, and Section~\ref{sec:experiments}
uses it precisely because its behaviour, unlike the policy's, is fully
explainable.

Yet the drones' autonomy is strictly \emph{executive and operational}, not
motivational. They choose their actions; they do not choose, weigh, reformulate
or abandon their goal, which was fixed by the designers, half at design time in
the reward function and half at training time in its distillation into values,
with no runtime process able to revise it. In delegation terms
\citep{dod2012autonomy}: the human specifies \emph{what} (find the target) and
every constraint on \emph{how} (the reward's economy of time, collision and
coverage); the system autonomously resolves everything below that line and
nothing above it. Two boundary facts make the point concrete. The drones cannot
refuse the mission or trade it off against anything, there being nothing else in
their value system; and the capability that looks most like initiative,
re-planning the division of labour when the team shrinks, is itself
goal-subordinate adaptation triggered by a belief update, in service of the same
fixed objective.

This placement, with high operational autonomy, zero motivational autonomy, no
runtime human in the loop and an opaque learned decision function, is exactly the
profile for which accountability questions are sharpest: the system's behaviour
is neither predictable enough to certify like an automatic device, nor deliberate
enough to interrogate like a reasoning agent. Section~\ref{sec:ethics} takes up
precisely this configuration.

\subsection{Learning in the multi-agent system}
\label{subsec:mas-learning}

Parameter sharing \citep{gupta2017cooperative} places this design at a specific
point of the CTDE continuum \citep{kraemer2016multi,lowe2017multi}: training-time
centralization consists \emph{only} of pooling experience into one replay buffer
and one set of weights, with no centralized critic, no access to teammates'
observations or actions, and no global state anywhere in the loss. Every
transition that trains the network is something some individual drone actually
experienced, and execution is correspondingly clean, since the batched action
selection is a wall-clock optimization with no informational content. The
deployment artefact is a fully decentralized policy in the strictest sense. The
same choice disposes of two MARL classics at once: the \emph{symmetry} of a
shared policy at a shared spawn is broken by the identity entry of the context
vector, and \emph{scalability} follows from never materializing the joint
anything, since sequential execution plus per-drone argmax means the joint action
space whose exponential growth motivates much of the literature simply never
appears, team size entering as a conditioning scalar rather than a
dimensionality. What sharing does \emph{not} remove is the moving-target problem,
each drone's optimal behaviour still depending on teammates' evolving behaviour,
and here the design accepts the standard empirical wager of parameter-shared
MARL: that co-evolution of one policy converges more gracefully than co-evolution
of $\nagents$.

\paragraph{The platform as a multi-agent laboratory.}
A final identity shift is worth recording: the same artefact that trains the
policy is, unchanged, an agent-based simulation platform
\citep{bonabeau2002agent}, a micro-level model whose macro-behaviour is studied
by controlled manipulation. The seeded, parameterized evaluation harness is
methodologically identical to a simulation experiment design: hold the population
of instances fixed, vary one factor ($\pfault$, $\rcomm$, team size), and measure
macro-level responses. In this reading Section~\ref{sec:experiments} is an
agent-based ``what-if'' study of the trained collective: what happens to a
searching society as its members become mortal.

\subsection{What a richer multi-agent design would add}
\label{subsec:mas-beyond}

The system's multi-agent machinery is deliberately minimal, and the analysis
above should not read as a claim that minimal is optimal. It is worth stating
what each classical layer this design omits would offer, and what it would cost,
both as intellectual honesty and as a roadmap.

\emph{Explicit communication protocols} \citep{finin1994kqml} would let drones
exchange intentions rather than facts (``I will sweep the north-east quadrant''),
enabling explicit task allocation and with it verifiable non-overlap guarantees
instead of statistically-learned dispersion. The cost is the machinery this
design conspicuously avoids: message semantics, conversation state, commitment
handling, and the robustness of all of it to the death of either partner
mid-protocol, a non-trivial concern in precisely this problem.

\emph{Deliberative architectures} \citep{rao1995bdi} would make goals and plans
explicit and inspectable. The gain is not primarily performance but
\emph{explainability}: a belief--desire--intention drone answers ``why did you go
north?'' with an intention trace, where the present drone can answer only with a
Q-value, a difference whose ethical weight
Section~\ref{subsec:eth-responsibility} makes precise. The cost is the classical
one: hand-authored plan libraries in a domain whose main difficulty is exactly
what resists concise symbolic specification.

\emph{Organizational structure}, meaning roles, formations and an explicit
division of the map, would replace learned conventions with designed ones,
gaining predictability and losing the seamless re-adaptation that makes the fault
response work: a role structure must be re-formed when its occupants die, whereas
the learned convention degrades continuously with $\hat n_i$.

\emph{Programmable coordination media}
\citep{gelernter1985generative,omicini2001tuple} offer perhaps the most natural
upgrade path, since the system already has a coordination medium (the fused map)
with a fixed policy (max-merge). Making that medium programmable, with a fusion
rule that could age knowledge, prioritize target and crash information under
bandwidth limits, or maintain provenance, would relocate coordination
intelligence from the opaque policy into an inspectable artefact, the move the
artefact literature explicitly recommends \citep{omicini2008artifacts}.

The common thread is a single trade-off, and it is the honest summary of this
section: every engineered layer buys transparency, verifiability and designer
control, and pays in brittleness and specification burden; the learned-minimal
design buys adaptivity across regimes that would each demand their own
specification, and pays in opacity. This project sits knowingly at the learned
end. The experimental section measures what that position buys; the ethical
section confronts what it costs.
